\chapter{O Sacramento da Confirmação}

Com o Batismo e a Eucaristia, o Sacramento da Confirmação constitui o conjunto dos ``sacramentos da iniciação cristã'', cuja unidade deve ser salvaguardada. Por isso, é preciso explicar aos fiéis que a recepção deste sacramento é necessária para a plenitude da graça batismal: os batizados ``pelo sacramento da Confirmação são mais perfeitamente vinculados à Igreja, enriquecidos com uma força especial do Espírito Santo, e deste modo ficam mais estritamente obrigados a difundir e a defender a fé por palavras e obras, como verdadeiras testemunhas de Cristo''. O Catecismo da Igreja Católica expõe nos parágrafos 1285 a 1321 a doutrina sobre a Confirmação: a sua origem na economia da salvação, os seus sinais e rito, os seus efeitos, as condições para recebê-la, o seu ministro e a síntese doutrinal dos elementos essenciais.

\section{A Confirmação na economia da salvação (1285--1292)}

No Antigo Testamento, os profetas anunciaram que o Espírito do Senhor repousaria sobre o Messias esperado, em vista da sua missão salvífica. Isaías descreve o Servo do Senhor sobre o qual repousa o Espírito --- Espírito de sabedoria e inteligência, Espírito de conselho e fortaleza, Espírito de ciência e temor do Senhor --- e o próprio Jesus aplica a si mesmo, na sinagoga de Nazaré, o oráculo de Isaías 61: ``O Espírito do Senhor está sobre mim, porque me ungiu para evangelizar os pobres''. A descida do Espírito Santo sobre Jesus, aquando do seu Batismo por João, foi o sinal de que era Ele o Messias, o Filho de Deus. Concebido pelo poder do Espírito Santo, toda a sua vida e toda a sua missão realizam-se numa comunhão total com o mesmo Espírito, que o Pai Lhe dá ``sem medida''.

Esta plenitude do Espírito, porém, não devia permanecer unicamente no Messias: devia ser comunicada a todo o povo messiânico. Repetidas vezes, Cristo prometeu esta efusão do Espírito, promessa que cumpriu, primeiro no dia de Páscoa, soprando sobre os Apóstolos, e depois, de modo mais esplêndido, no dia de Pentecostes. Cheios do Espírito Santo, os Apóstolos começaram a proclamar ``as maravilhas de Deus'', e Pedro declarou que esta efusão era o sinal dos tempos messiânicos, o cumprimento da promessa feita pelos profetas. Aqueles que então acreditaram na pregação apostólica e se fizeram batizar receberam, por sua vez, o dom do Espírito Santo.

A partir de então, os Apóstolos, para cumprir a vontade de Cristo, comunicavam aos neófitos, pela imposição das mãos, o dom do Espírito para completar a graça do Batismo. É por isso que, na Epístola aos Hebreus, se menciona, entre os elementos da primeira instrução cristã, a doutrina sobre os batismos e também sobre a imposição das mãos. A imposição das mãos é justificadamente reconhecida, pela Tradição católica, como a origem do sacramento da Confirmação, que ``de certo modo perpétua na Igreja a graça do Pentecostes''. Bem cedo, para melhor significar o dom do Espírito Santo, se acrescentou à imposição das mãos uma unção com óleo perfumado --- o crisma. Esta unção ilustra o nome de ``cristão'', que significa ``ungido'', indo buscar a sua origem ao próprio nome de Cristo, aquele que ``Deus ungiu com o Espírito Santo''.

Nos primeiros séculos, a Confirmação constituía geralmente uma única celebração com o Batismo, formando, segundo a expressão de São Cipriano, um ``sacramento duplo''. A multiplicação dos batismos de crianças em qualquer tempo do ano e o crescimento das paróquias rurais deixaram de permitir a presença do bispo em todas as celebrações batismais. No Ocidente, porque se desejava reservar ao bispo o completar do Batismo, instaurou-se a separação temporal dos dois sacramentos. O Oriente conservou os dois sacramentos unidos, de tal modo que a Confirmação é dada pelo sacerdote que batiza, mas com o \emph{myron} consagrado pelo bispo, exprimindo assim a unidade apostólica da Igreja. A prática das Igrejas do Oriente sublinha mais a unidade da iniciação cristã; a da Igreja latina exprime, com maior nitidez, a comunhão do novo cristão com o seu bispo --- garante e servidor da unidade, catolicidade e apostolicidade da Igreja ---, e assim a ligação com as origens apostólicas da Igreja de Cristo.

\section{Os sinais e o rito da Confirmação (1293--1301)}

No rito da Confirmação, convém considerar o sinal da unção e o que essa unção designa e imprime: o selo espiritual. A unção, na simbologia bíblica e antiga, é rica de numerosas significações: o óleo é sinal de abundância e de alegria, purifica e torna ágil; é sinal de cura, pois suaviza as contusões e as feridas; e torna radiante de beleza, saúde e força. Todos estes significados da unção com óleo se reencontram na vida sacramental: a unção antes do Batismo, com o óleo dos catecúmenos, significa purificação e fortalecimento; a unção dos enfermos exprime cura e conforto; a unção com o santo crisma depois do Batismo, na Confirmação e na Ordenação, é sinal duma consagração. Pela Confirmação, os cristãos --- os que são ungidos --- participam mais plenamente na missão de Jesus Cristo e na plenitude do Espírito, a fim de que toda a sua vida espalhe ``o bom odor de Cristo''.

Por esta unção, o confirmando recebe ``a marca'', o selo do Espírito Santo. O selo é o símbolo da pessoa, sinal de sua autoridade e de sua propriedade sobre um objeto; era assim que se marcavam os soldados com o selo do seu chefe e os escravos com o do seu dono. O próprio Cristo se declara marcado com o selo do Pai, e o cristão também é marcado: ``Foi Deus que nos concedeu a unção, nos marcou com o seu selo e depôs as arras do Espírito em nossos corações''. Este selo do Espírito Santo marca a pertença total a Cristo, a entrega para sempre ao seu serviço, e a promessa da proteção divina na grande prova escatológica.

Um momento importante que precede a celebração da Confirmação, mas que de certo modo faz parte dela, é a consagração do santo crisma. É o bispo que, na Quinta-Feira Santa, no decorrer da Missa Crismal, consagra o santo crisma para toda a sua diocese. Quando a Confirmação é celebrada separadamente do Batismo, como acontece no rito romano, a liturgia do sacramento começa pela renovação das promessas do Batismo e pela profissão de fé dos confirmandos, evidenciando claramente que a Confirmação se situa na continuação do Batismo. No caso do Batismo dum adulto, este recebe imediatamente a Confirmação e participa na Eucaristia.

No rito romano, o bispo estende as mãos sobre o grupo dos confirmandos, gesto que, desde o tempo dos Apóstolos, é sinal do dom do Espírito, invocando: ``Deus todo-poderoso, Pai de nosso Senhor Jesus Cristo, que, pela água e pelo Espírito Santo, destes uma vida nova a estes vossos servos e os libertastes do pecado, enviai sobre eles o Espírito Santo Paráclito; dai-lhes o espírito de sabedoria e de inteligência, o espírito de conselho e de fortaleza, o espírito de ciência e de piedade, e enchei-os do espírito do vosso temor''. O rito essencial consiste na unção do santo crisma sobre a fronte, feita com a imposição da mão, com as palavras: ``Recebe por este sinal o Espírito Santo, o Dom de Deus''. O ósculo da paz, com que termina o rito, significa e manifesta a comunhão eclesial com o bispo e com todos os fiéis.

\section{Os efeitos da Confirmação (1302--1305)}

O efeito do sacramento da Confirmação é uma efusão especial do Espírito Santo, tal como outrora foi concedida aos Apóstolos no dia de Pentecostes. Por isso, a Confirmação proporciona crescimento e aprofundamento da graça batismal: enraíza-nos mais profundamente na filiação divina, que nos leva a dizer ``Abbá, Pai!''; une-nos mais firmemente a Cristo; aumenta em nós os dons do Espírito Santo; torna mais perfeito o laço que nos une à Igreja; e dá-nos uma força especial do Espírito Santo para propagarmos e defendermos a fé pela palavra e pela obra, como verdadeiras testemunhas de Cristo, para confessarmos com valentia o nome de Cristo e para nunca nos envergonharmos da cruz. Santo Ambrósio resumiu esta doutrina: ``Lembra-te, pois, de que recebeste o sinal espiritual, o espírito de sabedoria e de entendimento, o espírito de conselho e de fortaleza, o espírito de ciência e de piedade, o espírito do santo temor, e guarda o que recebeste. Deus Pai marcou-te com o seu sinal, o Senhor Jesus Cristo confirmou-te e pôs no teu coração o penhor do Espírito''.

Tal como o Batismo, de que é a consumação, a Confirmação é dada uma só vez, pois imprime na alma uma marca espiritual indelével --- o ``caráter'' ---, que é sinal de que Jesus Cristo marcou o cristão com o selo do seu Espírito, revestindo-o da fortaleza do Alto para que seja sua testemunha. O caráter da Confirmação aperfeiçoa o sacerdócio comum dos fiéis, recebido no Batismo: ``o confirmado recebe a força de confessar a fé de Cristo publicamente, e como em virtude dum encargo oficial''. Por isso, a Confirmação não é uma mera formalidade de passagem, nem uma simples ``maturidade religiosa'', mas uma verdadeira efusão do Espírito que configura o fiel com Cristo Profeta, Sacerdote e Rei, e o envia em missão no coração do mundo.

\section{Quem pode receber a Confirmação (1306--1311)}

Todo o batizado ainda não confirmado pode e deve receber o sacramento da Confirmação. Uma vez que Batismo, Confirmação e Eucaristia formam uma unidade de iniciação, os fiéis têm obrigação de recebê-lo no tempo devido, pois, sem a Confirmação e a Eucaristia, o sacramento do Batismo é, sem dúvida, válido e eficaz, mas a iniciação cristã fica incompleta. O costume latino, desde há séculos, aponta ``a idade da discrição'' como referência para se receber a Confirmação; em perigo de morte, porém, devem confirmar-se as crianças, mesmo que ainda não tenham atingido essa idade.

Se por vezes se fala da Confirmação como ``sacramento da maturidade cristã'', não deve, no entanto, confundir-se a idade adulta da fé com a idade adulta do crescimento natural, nem esquecer-se que a graça batismal é uma graça de eleição gratuita e imerecida. São Tomás recorda isso: ``A idade do corpo não constitui um prejuízo para a alma. Por isso, mesmo na infância, o homem pode receber a perfeição da idade espiritual. E foi assim que muitas crianças, graças à fortaleza do Espírito Santo que tinham recebido, lutaram corajosamente e até ao sangue por Cristo''. A preparação para a Confirmação deve ter por fim conduzir o cristão a uma união mais íntima com Cristo e a uma familiaridade mais viva com o Espírito Santo, com a sua ação, os seus dons e os seus apelos, para melhor assumir as responsabilidades apostólicas. A catequese da Confirmação deve esforçar-se por despertar o sentido de pertença à Igreja, tanto universal como paroquial.

Para receber a Confirmação, é preciso estar em estado de graça. Convém recorrer ao sacramento da Penitência para ser purificado, em vista do dom do Espírito Santo, e uma oração mais intensa deve preparar para receber com docilidade e disponibilidade a força e as graças do Espírito Santo. Assim como no Batismo, convém que os candidatos à Confirmação procurem a ajuda espiritual dum padrinho ou duma madrinha. É conveniente que seja o mesmo do Batismo, para marcar bem a unidade dos dois sacramentos; o padrinho da Confirmação, a exemplo do padrinho do Batismo, exerce um verdadeiro ofício eclesial de acompanhamento espiritual.

\section{O ministro da Confirmação (1312--1314)}

O ministro originário da Confirmação é o bispo. No Oriente, é ordinariamente o sacerdote que batiza quem imediatamente confere a Confirmação, numa só e mesma celebração. Fá-lo, no entanto, com o santo crisma consagrado pelo patriarca ou pelo bispo, o que exprime a unidade apostólica da Igreja. Na Igreja latina aplica-se a mesma disciplina nos batismos de adultos ou quando é admitido à plena comunhão com a Igreja um batizado de outra comunidade cristã que não tenha recebido validamente o sacramento da Confirmação.

No rito latino, o ministro ordinário da Confirmação é o bispo, embora possa, em caso de necessidade, conceder a presbíteros a faculdade de administrá-la. É conveniente que seja ele mesmo a conferi-la, pois é por esse motivo que a celebração da Confirmação foi separada, no tempo, da do Batismo. Os bispos são os sucessores dos Apóstolos e receberam a plenitude do sacramento da Ordem; a administração deste sacramento por eles realça que ele tem como efeito unir mais estreitamente aqueles que o recebem à Igreja, às suas origens apostólicas e à sua missão de dar testemunho de Cristo. Se um cristão estiver em perigo de morte, qualquer sacerdote pode conferir-lhe a Confirmação, pois é vontade da Igreja que nenhum dos seus filhos, mesmo pequenino, parta deste mundo sem ter sido levado à perfeição pelo Espírito Santo com o dom da plenitude de Cristo.

\section{Em síntese (1315--1321)}

A Confirmação completa a graça batismal; ela é o sacramento que dá o Espírito Santo, para nos enraizar mais profundamente na filiação divina, incorporar-nos mais solidamente em Cristo, tornar mais firme o laço que nos prende à Igreja, associar-nos mais à sua missão e ajudar-nos a dar testemunho da fé cristã pela palavra acompanhada de obras. O fundamento escriturístico desta transmissão do Espírito está descrito no livro dos Atos dos Apóstolos: ``Quando os Apóstolos que estavam em Jerusalém ouviram dizer que a Samaria recebera a Palavra de Deus, enviaram-lhe Pedro e João. Quando chegaram lá, rezaram pelos samaritanos para que recebessem o Espírito Santo, que ainda não tinha descido sobre eles. Apenas tinham sido batizados em nome do Senhor Jesus. Então impunham-lhes as mãos e eles recebiam o Espírito Santo''.

A Confirmação, tal como o Batismo, imprime na alma do cristão um sinal espiritual --- o caráter indelével ---; por isso, só se pode receber este sacramento uma vez na vida. No Oriente, este sacramento é administrado imediatamente a seguir ao Batismo e é seguido da participação na Eucaristia; esta tradição põe em relevo a unidade dos três sacramentos da iniciação cristã. Na Igreja latina, este sacramento é administrado quando se atinge a idade da razão e, ordinariamente, a sua celebração é reservada ao bispo, significando que este sacramento vem robustecer o vínculo eclesial.

O candidato à Confirmação que atingiu a idade da razão deve professar a fé, estar em estado de graça, ter a intenção de receber o sacramento e estar preparado para assumir o seu papel de discípulo e testemunha de Cristo, na comunidade eclesial e nos assuntos temporais. O rito essencial da Confirmação é a unção com o santo crisma na fronte do batizado, com a imposição da mão do ministro e as palavras: ``Recebe por este sinal o Espírito Santo, o Dom de Deus''. Quando a Confirmação é celebrada separadamente do Batismo, a sua ligação com este sacramento é expressa, entre outras coisas, pela renovação dos compromissos batismais. A celebração da Confirmação no decorrer da Eucaristia sublinha a unidade perfeita dos três sacramentos da iniciação cristã, que juntos constituem a plena inserção do fiel no mistério de Cristo e na comunhão viva da Igreja.
